\documentclass[10pt]{extarticle}
\usepackage[a4paper,margin=1.5cm]{geometry}
\usepackage[T1]{fontenc}
\usepackage{amsmath,amssymb}
\usepackage{graphicx}
\usepackage[export]{adjustbox}
\usepackage{multicol}
\usepackage{enumitem}
\usepackage{booktabs,array,tabularx}
\setlength{\columnsep}{0.7cm}
\setlength{\parindent}{0pt}
\setlength{\parskip}{3pt}
\pagestyle{empty}
\makeatletter
% figure registry (figdefs.tex is generated by the extract step)
\newcommand{\qfigdef}[4]{\@namedef{qf@file@#1}{#2}\@namedef{qf@w@#1}{#3}\@namedef{qf@wide@#1}{#4}}
\newcommand{\qf@check}[1]{\@ifundefined{qf@file@#1}{\PackageError{score2}{Unknown figure id #1}{}}{}}
\newcommand{\qf@img}[4][]{\includegraphics[width=#2\dimexpr\@nameuse{qf@w@#4}pt\relax,max width=#3,max height=0.45\textheight,#1]{\@nameuse{qf@file@#4}}}
% Full-width questions: finalize wraps every question that contains a wide figure in \qfullbegin...\qfullend,
% so the WHOLE question (stem, figures, options) is set across the page in reading order, then the two columns
% resume. The question's head - its number, stem and first figure (\qfullsplit marks the end of the head) -
% is kept on one page, so a number is never left at a page foot with its figure overleaf; the rest may break
% across pages normally, so no big gaps are left. Without a \qfullsplit, the whole question is kept together
% if it is under 45% of a page. Never write these yourself.
\newif\ifqf@full
\newif\ifqf@boxed
\newsavebox{\qf@box}
\newcommand{\qf@open}{\global\qf@fulltrue\par\global\qf@boxedtrue
  \setbox\qf@box\vbox\bgroup\hsize\textwidth\linewidth\textwidth}
\newcommand{\qf@emit}[1]{\par\egroup\global\qf@boxedfalse
  \ifdim\dimexpr\ht\qf@box+\dp\qf@box\relax<#1\textheight\noindent\usebox{\qf@box}\par\else\unvbox\qf@box\fi}
\newcommand{\qfullsplit}{\ifqf@boxed\qf@emit{0.85}\fi}
\newcommand{\qf@close}{\ifqf@boxed\qf@emit{0.45}\fi\global\qf@fullfalse}
\newcommand{\qfullbegin}{\end{multicols}\qf@open}
\newcommand{\qfullend}{\qf@close\begin{multicols}{2}}
\newcommand{\qfullstart}{\qf@open}
\newcommand{\qfulldone}{\qf@close}
% \qfig[scale]{id}: block figure, centred, directly where it is written (size set by the tool).
% Fallback if a wide figure is not inside a full-width question: break out of the columns in place - no floats.
\newcommand{\qfig}[2][1]{\qf@check{#2}%
  \ifqf@full\qf@blk{#1}{#2}\else
  \ifnum\@nameuse{qf@wide@#2}=1 \expandafter\@firstoftwo\else\expandafter\@secondoftwo\fi
  {\end{multicols}\begin{center}\qf@img{#1}{\textwidth}{#2}\end{center}\begin{multicols}{2}}%
  {\qf@blk{#1}{#2}}\fi}
\newcommand{\qf@blk}[2]{\par\smallskip{\centering\qf@img{#1}{\linewidth}{#2}\par}\smallskip}
% \qfigi[scale]{id}: inline figure (inside option lists / tables / minipages); never breaks columns.
\newcommand{\qfigi}[2][1]{\qf@check{#2}\qf@img[valign=c]{#1}{\linewidth}{#2}}
\makeatother
% structure macros - lint reads these to check completeness and figure placement
\newcommand{\papersec}[1]{\par\bigskip{\centering\Large\bfseries PAPER-#1\par}\medskip}
\newcommand{\partsec}[2]{\par\medskip\noindent{\large\bfseries PART-#1: #2}\par\smallskip}
\newcommand{\sect}[1]{\par\smallskip\noindent\textbf{#1}\par}
\newcommand{\qn}[1]{\par\smallskip\noindent\textbf{#1.}~}
\newcommand{\paraq}[2]{\par\smallskip\noindent\textbf{Paragraph for Questions #1 and #2}\par}
\newenvironment{qoptions}{\begin{itemize}[leftmargin=1.8em,itemsep=1pt,topsep=2pt,label={}]}{\end{itemize}}
\newcommand{\opt}[2]{\item[(#1)] #2}
% figure registry (generated by score2_tools.py finalize)
\qfigdef{p1_phy_s1_q01_1}{TMAS - Test 01 (09 Jul 2023)-img/p1_phy_s1_q01_1.png}{55.6}{0}
\qfigdef{p1_phy_s1_q03_1}{TMAS - Test 01 (09 Jul 2023)-img/p1_phy_s1_q03_1.png}{63.4}{0}
\qfigdef{p1_phy_s1_q04_1}{TMAS - Test 01 (09 Jul 2023)-img/p1_phy_s1_q04_1.png}{166.3}{0}
\qfigdef{p1_phy_s1_q05_1}{TMAS - Test 01 (09 Jul 2023)-img/p1_phy_s1_q05_1.png}{137.6}{0}
\qfigdef{p1_phy_s1_q06_1}{TMAS - Test 01 (09 Jul 2023)-img/p1_phy_s1_q06_1.png}{165.3}{0}
\qfigdef{p1_phy_s1_q07_1}{TMAS - Test 01 (09 Jul 2023)-img/p1_phy_s1_q07_1.png}{122.0}{0}
\qfigdef{p1_phy_s1_q08_1}{TMAS - Test 01 (09 Jul 2023)-img/p1_phy_s1_q08_1.png}{145.0}{0}
\qfigdef{p1_phy_s1_q09_1}{TMAS - Test 01 (09 Jul 2023)-img/p1_phy_s1_q09_1.png}{80.1}{0}
\qfigdef{p1_phy_s1_q10_1}{TMAS - Test 01 (09 Jul 2023)-img/p1_phy_s1_q10_1.png}{125.6}{0}
\qfigdef{p1_phy_s1_q11_1}{TMAS - Test 01 (09 Jul 2023)-img/p1_phy_s1_q11_1.png}{91.1}{0}
\qfigdef{p1_phy_s1_q12_1}{TMAS - Test 01 (09 Jul 2023)-img/p1_phy_s1_q12_1.png}{88.0}{0}
\qfigdef{p1_phy_s1_q13_1}{TMAS - Test 01 (09 Jul 2023)-img/p1_phy_s1_q13_1.png}{90.3}{0}
\qfigdef{p1_phy_s1_q17_1}{TMAS - Test 01 (09 Jul 2023)-img/p1_phy_s1_q17_1.png}{87.9}{0}
\qfigdef{p1_phy_s1_q18_1}{TMAS - Test 01 (09 Jul 2023)-img/p1_phy_s1_q18_1.png}{103.1}{0}
\qfigdef{p1_phy_s2_q01_1}{TMAS - Test 01 (09 Jul 2023)-img/p1_phy_s2_q01_1.png}{86.9}{0}
\qfigdef{p1_phy_s2_q04_1}{TMAS - Test 01 (09 Jul 2023)-img/p1_phy_s2_q04_1.png}{145.7}{0}
\qfigdef{p1_phy_s2_q08_1}{TMAS - Test 01 (09 Jul 2023)-img/p1_phy_s2_q08_1.png}{122.9}{0}
\qfigdef{p1_phy_s2_q10_1}{TMAS - Test 01 (09 Jul 2023)-img/p1_phy_s2_q10_1.png}{92.6}{0}
\qfigdef{p1_chem_s1_q01_1}{TMAS - Test 01 (09 Jul 2023)-img/p1_chem_s1_q01_1.png}{167.0}{0}
\qfigdef{p1_chem_s1_q14_1}{TMAS - Test 01 (09 Jul 2023)-img/p1_chem_s1_q14_1.png}{43.3}{0}
\qfigdef{p1_chem_s1_q14_2}{TMAS - Test 01 (09 Jul 2023)-img/p1_chem_s1_q14_2.png}{39.1}{0}
\qfigdef{p1_chem_s1_q14_3}{TMAS - Test 01 (09 Jul 2023)-img/p1_chem_s1_q14_3.png}{41.7}{0}
\qfigdef{p1_chem_s1_q14_4}{TMAS - Test 01 (09 Jul 2023)-img/p1_chem_s1_q14_4.png}{42.7}{0}
\qfigdef{p1_chem_s1_q15_1}{TMAS - Test 01 (09 Jul 2023)-img/p1_chem_s1_q15_1.png}{165.1}{0}
\qfigdef{p1_chem_s1_q16_1}{TMAS - Test 01 (09 Jul 2023)-img/p1_chem_s1_q16_1.png}{88.9}{0}
\qfigdef{p1_chem_s1_q17_1}{TMAS - Test 01 (09 Jul 2023)-img/p1_chem_s1_q17_1.png}{116.6}{0}
\qfigdef{p1_chem_s1_q17_A}{TMAS - Test 01 (09 Jul 2023)-img/p1_chem_s1_q17_A.png}{39.7}{0}
\qfigdef{p1_chem_s1_q17_B}{TMAS - Test 01 (09 Jul 2023)-img/p1_chem_s1_q17_B.png}{26.2}{0}
\qfigdef{p1_chem_s1_q17_C}{TMAS - Test 01 (09 Jul 2023)-img/p1_chem_s1_q17_C.png}{36.0}{0}
\qfigdef{p1_chem_s1_q17_D}{TMAS - Test 01 (09 Jul 2023)-img/p1_chem_s1_q17_D.png}{33.1}{0}
\qfigdef{p1_chem_s1_q18_1}{TMAS - Test 01 (09 Jul 2023)-img/p1_chem_s1_q18_1.png}{148.0}{0}
\qfigdef{p1_chem_s1_q19_A}{TMAS - Test 01 (09 Jul 2023)-img/p1_chem_s1_q19_A.png}{95.4}{0}
\qfigdef{p1_chem_s1_q19_3}{TMAS - Test 01 (09 Jul 2023)-img/p1_chem_s1_q19_3.png}{59.5}{0}
\qfigdef{p1_chem_s1_q19_B}{TMAS - Test 01 (09 Jul 2023)-img/p1_chem_s1_q19_B.png}{142.5}{0}
\qfigdef{p1_chem_s1_q19_C}{TMAS - Test 01 (09 Jul 2023)-img/p1_chem_s1_q19_C.png}{65.8}{0}
\qfigdef{p1_chem_s1_q19_6}{TMAS - Test 01 (09 Jul 2023)-img/p1_chem_s1_q19_6.png}{37.7}{0}
\qfigdef{p1_chem_s1_q19_D}{TMAS - Test 01 (09 Jul 2023)-img/p1_chem_s1_q19_D.png}{93.8}{0}
\qfigdef{p1_chem_s1_q20_1}{TMAS - Test 01 (09 Jul 2023)-img/p1_chem_s1_q20_1.png}{92.8}{0}
\qfigdef{p1_chem_s1_q20_2}{TMAS - Test 01 (09 Jul 2023)-img/p1_chem_s1_q20_2.png}{166.5}{0}
\qfigdef{p1_chem_s1_q20_A}{TMAS - Test 01 (09 Jul 2023)-img/p1_chem_s1_q20_A.png}{47.2}{0}
\qfigdef{p1_chem_s1_q20_B}{TMAS - Test 01 (09 Jul 2023)-img/p1_chem_s1_q20_B.png}{48.9}{0}
\qfigdef{p1_chem_s1_q20_C}{TMAS - Test 01 (09 Jul 2023)-img/p1_chem_s1_q20_C.png}{72.7}{0}
\qfigdef{p1_chem_s1_q20_D}{TMAS - Test 01 (09 Jul 2023)-img/p1_chem_s1_q20_D.png}{48.9}{0}
\qfigdef{p1_chem_s2_q03_1}{TMAS - Test 01 (09 Jul 2023)-img/p1_chem_s2_q03_1.png}{114.8}{0}
\qfigdef{p1_chem_s2_q08_1}{TMAS - Test 01 (09 Jul 2023)-img/p1_chem_s2_q08_1.png}{138.1}{0}
\qfigdef{p1_chem_s2_q10_1}{TMAS - Test 01 (09 Jul 2023)-img/p1_chem_s2_q10_1.png}{37.2}{0}
\qfigdef{p1_math_s1_q10_1}{TMAS - Test 01 (09 Jul 2023)-img/p1_math_s1_q10_1.png}{85.7}{0}
\qfigdef{p2_phy_s1_q01_1}{TMAS - Test 01 (09 Jul 2023)-img/p2_phy_s1_q01_1.png}{184.1}{0}
\qfigdef{p2_phy_s1_q01_C}{TMAS - Test 01 (09 Jul 2023)-img/p2_phy_s1_q01_C.png}{103.9}{0}
\qfigdef{p2_phy_s1_q01_D}{TMAS - Test 01 (09 Jul 2023)-img/p2_phy_s1_q01_D.png}{104.4}{0}
\qfigdef{p2_phy_s1_q02_1}{TMAS - Test 01 (09 Jul 2023)-img/p2_phy_s1_q02_1.png}{126.9}{0}
\qfigdef{p2_phy_s1_q03_1}{TMAS - Test 01 (09 Jul 2023)-img/p2_phy_s1_q03_1.png}{129.2}{0}
\qfigdef{p2_phy_s1_q06_1}{TMAS - Test 01 (09 Jul 2023)-img/p2_phy_s1_q06_1.png}{105.3}{0}
\qfigdef{p2_phy_s1_q07_1}{TMAS - Test 01 (09 Jul 2023)-img/p2_phy_s1_q07_1.png}{108.3}{0}
\qfigdef{p2_phy_s1_q07_2}{TMAS - Test 01 (09 Jul 2023)-img/p2_phy_s1_q07_2.png}{134.9}{0}
\qfigdef{p2_phy_s2_para8_10_1}{TMAS - Test 01 (09 Jul 2023)-img/p2_phy_s2_para8_10_1.png}{80.4}{0}
\qfigdef{p2_phy_s2_para8_10_2}{TMAS - Test 01 (09 Jul 2023)-img/p2_phy_s2_para8_10_2.png}{56.9}{0}
\qfigdef{p2_phy_s2_para8_10_3}{TMAS - Test 01 (09 Jul 2023)-img/p2_phy_s2_para8_10_3.png}{92.2}{0}
\qfigdef{p2_phy_s2_para8_10_4}{TMAS - Test 01 (09 Jul 2023)-img/p2_phy_s2_para8_10_4.png}{134.7}{0}
\qfigdef{p2_phy_s3_q01_1}{TMAS - Test 01 (09 Jul 2023)-img/p2_phy_s3_q01_1.png}{133.0}{0}
\qfigdef{p2_phy_s3_q03_1}{TMAS - Test 01 (09 Jul 2023)-img/p2_phy_s3_q03_1.png}{112.4}{0}
\qfigdef{p2_phy_s3_q04_1}{TMAS - Test 01 (09 Jul 2023)-img/p2_phy_s3_q04_1.png}{107.3}{0}
\qfigdef{p2_phy_s3_q05_1}{TMAS - Test 01 (09 Jul 2023)-img/p2_phy_s3_q05_1.png}{69.3}{0}
\qfigdef{p2_chem_s1_q01_A}{TMAS - Test 01 (09 Jul 2023)-img/p2_chem_s1_q01_A.png}{157.9}{0}
\qfigdef{p2_chem_s1_q01_B}{TMAS - Test 01 (09 Jul 2023)-img/p2_chem_s1_q01_B.png}{105.8}{0}
\qfigdef{p2_chem_s1_q01_C}{TMAS - Test 01 (09 Jul 2023)-img/p2_chem_s1_q01_C.png}{128.8}{0}
\qfigdef{p2_chem_s1_q01_D}{TMAS - Test 01 (09 Jul 2023)-img/p2_chem_s1_q01_D.png}{130.3}{0}
\qfigdef{p2_chem_s1_q02_A}{TMAS - Test 01 (09 Jul 2023)-img/p2_chem_s1_q02_A.png}{89.9}{0}
\qfigdef{p2_chem_s1_q02_B}{TMAS - Test 01 (09 Jul 2023)-img/p2_chem_s1_q02_B.png}{128.2}{0}
\qfigdef{p2_chem_s1_q02_C}{TMAS - Test 01 (09 Jul 2023)-img/p2_chem_s1_q02_C.png}{158.4}{0}
\qfigdef{p2_chem_s1_q02_D}{TMAS - Test 01 (09 Jul 2023)-img/p2_chem_s1_q02_D.png}{190.0}{0}
\qfigdef{p2_chem_s1_q04_A}{TMAS - Test 01 (09 Jul 2023)-img/p2_chem_s1_q04_A.png}{72.0}{0}
\qfigdef{p2_chem_s1_q04_B}{TMAS - Test 01 (09 Jul 2023)-img/p2_chem_s1_q04_B.png}{119.3}{0}
\qfigdef{p2_chem_s1_q04_C}{TMAS - Test 01 (09 Jul 2023)-img/p2_chem_s1_q04_C.png}{74.5}{0}
\qfigdef{p2_chem_s1_q04_D}{TMAS - Test 01 (09 Jul 2023)-img/p2_chem_s1_q04_D.png}{97.5}{0}
\qfigdef{p2_chem_s1_q05_1}{TMAS - Test 01 (09 Jul 2023)-img/p2_chem_s1_q05_1.png}{116.5}{0}
\qfigdef{p2_chem_s3_q01_1}{TMAS - Test 01 (09 Jul 2023)-img/p2_chem_s3_q01_1.png}{95.6}{0}
\qfigdef{p2_chem_s3_q02_1}{TMAS - Test 01 (09 Jul 2023)-img/p2_chem_s3_q02_1.png}{136.1}{0}
\qfigdef{p2_chem_s3_q02_2}{TMAS - Test 01 (09 Jul 2023)-img/p2_chem_s3_q02_2.png}{121.3}{0}
\qfigdef{p2_chem_s3_q02_3}{TMAS - Test 01 (09 Jul 2023)-img/p2_chem_s3_q02_3.png}{95.6}{0}
\qfigdef{p2_chem_s3_q02_4}{TMAS - Test 01 (09 Jul 2023)-img/p2_chem_s3_q02_4.png}{206.7}{0}
\qfigdef{p2_chem_s3_q02_5}{TMAS - Test 01 (09 Jul 2023)-img/p2_chem_s3_q02_5.png}{87.7}{0}
\qfigdef{p2_chem_s3_q02_6}{TMAS - Test 01 (09 Jul 2023)-img/p2_chem_s3_q02_6.png}{184.1}{0}
\qfigdef{p2_chem_s3_q02_7}{TMAS - Test 01 (09 Jul 2023)-img/p2_chem_s3_q02_7.png}{199.6}{0}
\qfigdef{p2_chem_s3_q02_8}{TMAS - Test 01 (09 Jul 2023)-img/p2_chem_s3_q02_8.png}{197.0}{0}
\qfigdef{p2_chem_s3_q02_9}{TMAS - Test 01 (09 Jul 2023)-img/p2_chem_s3_q02_9.png}{161.1}{0}
\qfigdef{p2_chem_s3_q04_1}{TMAS - Test 01 (09 Jul 2023)-img/p2_chem_s3_q04_1.png}{143.5}{0}
\qfigdef{p2_chem_s3_q05_1}{TMAS - Test 01 (09 Jul 2023)-img/p2_chem_s3_q05_1.png}{44.9}{0}
\qfigdef{p2_chem_s3_q05_2}{TMAS - Test 01 (09 Jul 2023)-img/p2_chem_s3_q05_2.png}{47.8}{0}
\qfigdef{p2_chem_s3_q05_3}{TMAS - Test 01 (09 Jul 2023)-img/p2_chem_s3_q05_3.png}{71.5}{0}
\qfigdef{p2_chem_s3_q05_6}{TMAS - Test 01 (09 Jul 2023)-img/p2_chem_s3_q05_6.png}{42.4}{0}
\qfigdef{p2_chem_s3_q05_4}{TMAS - Test 01 (09 Jul 2023)-img/p2_chem_s3_q05_4.png}{68.9}{0}
\qfigdef{p2_chem_s3_q05_5}{TMAS - Test 01 (09 Jul 2023)-img/p2_chem_s3_q05_5.png}{67.6}{0}
\qfigdef{p2_chem_s3_q05_9}{TMAS - Test 01 (09 Jul 2023)-img/p2_chem_s3_q05_9.png}{46.8}{0}
\qfigdef{p2_chem_s3_q05_7}{TMAS - Test 01 (09 Jul 2023)-img/p2_chem_s3_q05_7.png}{61.0}{0}
\qfigdef{p2_chem_s3_q05_8}{TMAS - Test 01 (09 Jul 2023)-img/p2_chem_s3_q05_8.png}{62.2}{0}
\qfigdef{p2_chem_s3_q05_10}{TMAS - Test 01 (09 Jul 2023)-img/p2_chem_s3_q05_10.png}{67.6}{0}
\raggedcolumns
\makeatletter\@beginparpenalty=10000\makeatother
\newcolumntype{L}[1]{>{\raggedright\arraybackslash}p{#1}}
\begin{document}
\begin{center}{\Large\bfseries TMAS --- Test 01 (09 Jul 2023)}\end{center}
\begin{multicols}{2}
\raggedcolumns
\papersec{1}
\partsec{1}{PHYSICS}
\sect{SECTION-I}
\qn{1} Electric field at point $P$ is given by $\vec E=E_0r\,\hat r$. The total flux through the given cylinder of radius $R$ and height $h$ is :-
\qfig{p1_phy_s1_q01_1}
\begin{qoptions}
\opt{A}{$E_0\pi R^2h$}
\opt{B}{$2E_0\pi R^2h$}
\opt{C}{$3E_0\pi R^2h$}
\opt{D}{$4E_0\pi R^2h$}
\end{qoptions}
\qn{2} \textbf{Statement-1 :} Work done by electric field on moving a positive charge on an equipotential surface is always zero.\\
and\\
\textbf{Statement-2 :} Electric lines of forces are perpendicular to equipotential surface.
\begin{qoptions}
\opt{A}{Statement--1 is True, Statement--2 is True ; Statement--2 is a correct explanation for Statement--1.}
\opt{B}{Statement--1 is True, Statement--2 is True ; Statement--2 is not a correct explanation for Statement--1.}
\opt{C}{Statement--1 is True, Statement--2 is False.}
\opt{D}{Statement--1 is False, Statement--2 is True.}
\end{qoptions}
\qn{3} Four charges are arranged at the corners of a square $ABCD$, as shown. The net force on positive charge kept at the centre of the square is
\qfig{p1_phy_s1_q03_1}
\begin{qoptions}
\opt{A}{zero}
\opt{B}{along diagonal $AC$}
\opt{C}{along diagonal $BD$}
\opt{D}{perpendicular to the side $AB$}
\end{qoptions}
\qn{4} The density of charge at point $P$ on the conductor is $\sigma$. The resultant electric field near $P$ (just outside) will be :-
\qfig{p1_phy_s1_q04_1}
\begin{qoptions}
\opt{A}{$\dfrac{kQ_1}{r_1^2}\left(-\hat j\right)+\dfrac{kQ_2}{r_2^2}\left(-\hat i\right)$}
\opt{B}{$\dfrac{kQ_1}{r_1^2}\left(-\hat j\right)+\dfrac{kQ_2}{r_2^2}\left(-\hat i\right)+\dfrac\sigma{\in_0}\left(\hat i\right)$}
\opt{C}{$\dfrac\sigma{\in_0}\hat i$}
\opt{D}{$\dfrac{kQ_1}{r_1^2}\left(\hat j\right)+\dfrac{kQ_2}{r_2^2}\left(\hat i\right)+\dfrac\sigma{\in_0}\left(-\hat i\right)$}
\end{qoptions}
\qn{5} Find out electric flux passing from a closed Gaussian surface as shown in figure :
\qfig{p1_phy_s1_q05_1}
\begin{qoptions}
\opt{A}{$\dfrac{q_1+q_5-q_6}{\varepsilon_0}$}
\opt{B}{$\dfrac{q_1+q_5+q_6}{\varepsilon_0}$}
\opt{C}{$\dfrac{q_1+q_2+q_3+q_4+q_5+q_6}{\varepsilon_0}$}
\opt{D}{$\dfrac{q_1+q_2-q_3+q_4+q_5-q_6}{\varepsilon_0}$}
\end{qoptions}
\qn{6} Two point charges $Q_1$ and $Q_2$ lie along a line at a distance $a$ from each other. Figure shows the potential variation along line of charges then which of the following statement is incorrect?
\qfig{p1_phy_s1_q06_1}
\begin{qoptions}
\opt{A}{Electric field is zero at point 3}
\opt{B}{$Q_1$ is positive and $Q_2$ is negative}
\opt{C}{$Q_1$ is negative and $Q_2$ is positive}
\opt{D}{$|Q_1|>|Q_2|$}
\end{qoptions}
\qn{7} In the circuit shown in the figure, if the potential at point $A$ is taken to be $-2$V, the potential at point $B$ is :-
\qfig{p1_phy_s1_q07_1}
\begin{qoptions}
\opt{A}{$-3$V}
\opt{B}{$-7$V}
\opt{C}{$-4$V}
\opt{D}{None}
\end{qoptions}
\qn{8} Net resistance between $X$ and $Y$ is --
\qfig{p1_phy_s1_q08_1}
\begin{qoptions}
\opt{A}{5 $\Omega$}
\opt{B}{10 $\Omega$}
\opt{C}{15 $\Omega$}
\opt{D}{60 $\Omega$}
\end{qoptions}
\qn{9} In the given bridge the value of $X$ for which the potential difference between the point $B$ and $D$ will be zero is:
\qfig{p1_phy_s1_q09_1}
\begin{qoptions}
\opt{A}{5 $\Omega$}
\opt{B}{10 $\Omega$}
\opt{C}{15 $\Omega$}
\opt{D}{20 $\Omega$}
\end{qoptions}
\qn{10} A given resistance is connected with a battery of emf 170 volt exactly with negligible internal resistance. Find the value of current flowing through resistor :-
\qfig{p1_phy_s1_q10_1}
\begin{qoptions}
\opt{A}{1 mA $\pm$ 5\%}
\opt{B}{17 mA $\pm$ 5\%}
\opt{C}{2 mA $\pm$ 10\%}
\opt{D}{15 mA $\pm$ 10\%}
\end{qoptions}
\qn{11} A potentiometer has a resistance wire of 10 m length and resistance 10$\Omega$ as shown in figure. A circuit is setup as shown. For null deflection in galvanometer length $AP$ is :
\qfig{p1_phy_s1_q11_1}
\begin{qoptions}
\opt{A}{2m}
\opt{B}{4.5 m}
\opt{C}{6.5 m}
\opt{D}{None of these}
\end{qoptions}
\qn{12} For a wheatstone bridge shown; $R_1>R_2$ and $R_3=R_4$. The direction of current in wire $AB$ is:
\qfig{p1_phy_s1_q12_1}
\begin{qoptions}
\opt{A}{From $A$ to $B$}
\opt{B}{No current flows in $AB$}
\opt{C}{From $B$ to $A$}
\opt{D}{Data is not sufficient}
\end{qoptions}
\qn{13} The earth is moving on an elliptical path, whose one focus sun is situated as shown in figure. If $AS=r_{\min}$ and $SB=r_{\max}$, the sun and earth system obey the Kepler's law, the square of time-period is directly proportional to
\qfig{p1_phy_s1_q13_1}
\begin{qoptions}
\opt{A}{$(r_{\min})^3$}
\opt{B}{$(r_{\max})^3$}
\opt{C}{$\left(\dfrac{r_{\min}+r_{\max}}2\right)^3$}
\opt{D}{$\left(\dfrac{r_{\min}r_{\max}}{r_{\min}+r_{\max}}\right)^3$}
\end{qoptions}
\qn{14} When a satellite in a circular orbit around the earth enters the atmospheric region, it encounters small air resistance to its motion. Then
\begin{qoptions}
\opt{A}{It kinetic energy increases}
\opt{B}{It kinetic energy decreases}
\opt{C}{Its angular momentum about the earth increases}
\opt{D}{Its period of revolution around the earth increases}
\end{qoptions}
\qn{15} Two spherical stars each of mass $m$ with their centres at distance $D$ revolve under mutual gravitational attraction about the point midway between their centres. The period of revolution will be :-
\begin{qoptions}
\opt{A}{$\pi\sqrt{\dfrac{D^3}{GM}}$}
\opt{B}{$2\pi\sqrt{\dfrac{D^3}{2GM}}$}
\opt{C}{$2\pi\sqrt{\dfrac{2D^3}{GM}}$}
\opt{D}{$\pi\sqrt{\dfrac{GM}{D^3}}$}
\end{qoptions}
\qn{16} A satellite is revolving round the earth in an orbit of radius $r$ with time period $T$. If the satellite is revolving round the earth in an orbit of radius $r+\Delta r$ ($\Delta r\ll r$) with time period $T+\Delta T$ then :-
\begin{qoptions}
\opt{A}{$\dfrac{\Delta T}T=\dfrac32\dfrac{\Delta r}r$}
\opt{B}{$\dfrac{\Delta T}T=\dfrac23\dfrac{\Delta r}r$}
\opt{C}{$\dfrac{\Delta T}T=\dfrac{\Delta r}r$}
\opt{D}{$\dfrac{\Delta T}T=-\dfrac{\Delta r}r$}
\end{qoptions}
\qn{17} A planet revolves around the sun in an elliptical orbit. The linear speed of the planet will be maximum at
\qfig{p1_phy_s1_q17_1}
\begin{qoptions}
\opt{A}{$D$}
\opt{B}{$B$}
\opt{C}{$A$}
\opt{D}{$C$}
\end{qoptions}
\qn{18} Charge distribution on a ring in $x$-$y$ plane is shown in the figure. Then electric potential at $O$ is :
\qfig{p1_phy_s1_q18_1}
\begin{qoptions}
\opt{A}{$\dfrac Q{4\pi\in_0r}$}
\opt{B}{Zero}
\opt{C}{$\dfrac Q{2\pi\in_0r}$}
\opt{D}{$\dfrac{3Q}{4\pi\in_0r}$}
\end{qoptions}
\qn{19} The equation of an equipotential line in an electric field is $y=2x$, then the electric field strength vector at $(1,2)$ may be :
\begin{qoptions}
\opt{A}{$4\hat i+3\hat j$}
\opt{B}{$4\hat i+8\hat j$}
\opt{C}{$8\hat i+4\hat j$}
\opt{D}{$-8\hat i+4\hat j$}
\end{qoptions}
\qn{20} Two identical metal balls with charge $+2Q$ and $-Q$ are separated by some distance, and exert a force $F$ on each other. They are joined by a long conducting wire, which is then removed. The force between them will now be
\begin{qoptions}
\opt{A}{$F$}
\opt{B}{$F/2$}
\opt{C}{$F/4$}
\opt{D}{$F/8$}
\end{qoptions}
\sect{SECTION-II}
\qn{1} An electric field is given by $\vec E=5\hat i-(7y+8)\hat j-3\hat k$ (S.I. unit) crosses gaussian cube having side length 2m as shown in the figure. Then the net charge enclosed in the cube in $-x\in_0$ (S.I. unit) then find $x$.
\qfig{p1_phy_s2_q01_1}
\qn{2} Four $+q$ charges are fixed on corner of square of side $a$. A fifth charge $+Q$ is brought by external agent slowly from infinity and placed at centre of square. Work done by external agent is $\dfrac{\sqrt\alpha kQq}a$ (where $k=\dfrac1{4\pi\in_0}$). Fill $\alpha$.
\qn{3} Two charges $+q_1$ and $-q_2$ are placed at $A$ and $B$ respectively. A line of force emanates from $q_1$ at an angle $\alpha$ with the line $AB$. If it terminates at $-q_2$ at an angle $P\sin^{-1}\left[\sqrt{\dfrac{q_1}{q_2}}\sin\dfrac\alpha2\right]$. Find the value of $P$
\qn{4} An infinitely long wire with linear charge density $\lambda_1$ and a segment of length $\ell$ with charge density $\lambda_2$ lie in a plane at right angle to each other and seperated by a distance $r_0$ as shown in figure. The force with which these two interact is $\dfrac{\lambda_1\lambda_2}{4\pi\varepsilon_0}\,\ell n\,x$. If $\ell=2r_0$, then find the value of $x$.
\qfig{p1_phy_s2_q04_1}
\qn{5} Find out the electric flux (S.I. unit) through an area 10 m$^2$ lying in $XY$ plane due to an electric field $\vec E=2\hat i-10\hat j+5\hat k$ (S.I. unit).
\qn{6} The density of core of a planet is $\rho_1$ and that of the outer shell is $\rho_2$. The radii of the core and outer shell of the planet are $R$ and $2R$ respectively. If acceleration due to gravity at the surface is half as at a depth $R$ then find $\dfrac{\rho_1}{\rho_2}$.
\qn{7} The gravitational field in a region is given by $\vec E=(3\hat i-4\hat j)$ N/kg. Find out the work done (in joule) in displacing a particle by 1 m, along the line $4y=3x+9$ :
\qn{8} Consider the potentiometer circuit arranged as in the figure. The potentiometer wire is 600 cm long. Resistance of wire $AB$ is $15r$. At what distance (in cm) from Point $A$ should the jockey touch the wire to get zero deflection in the galvanometer.
\qfig{p1_phy_s2_q08_1}
\qn{9} The current in a wire varies with time according to the equation $I=4+2t$ where $I$ is in ampere and $t$ is in sec. Calculate the quantity of charge (in columb) which has passed through a cross-section of the wire during the time $t=2$ sec. to $t=6$ sec.
\qn{10} In the circuit shown in the adjoining figure, the reading of ideal ammeter $A$ is :
\qfig{p1_phy_s2_q10_1}
\partsec{2}{CHEMISTRY}
\sect{SECTION-I}
\qn{1} Select the most appropriate option about ligand considering below given reactions : (tu = thiourea)
\qfig{p1_chem_s1_q01_1}
\begin{qoptions}
\opt{A}{In reaction-1 product, thiosulphate act as bidentate ligand.}
\opt{B}{In reaction-2 product, thiosulphate act as bidentate ligand.}
\opt{C}{In reaction-1 product, tu act as bidentate ligand.}
\opt{D}{In reaction-2 product, tu act as bidentate ligand.}
\end{qoptions}
\qn{2} Correct statement is :\\
(i) In 13$^{\text{th}}$ group stability of lower oxidation state increases down the group\\
(ii) Pentavalent bismuth act as reducing agent\\
(iii) $\mathrm{Pb^{4+}}$ is good oxidising agent because $\mathrm{Pb^{4+}}$ have tendency to gain electron and convert into more stable $\mathrm{Pb^{2+}}$ due to inert pair effect.\\
(iv) $\mathrm{BiF_3}$ exist where as $\mathrm{BiF_5}$ not exist.
\begin{qoptions}
\opt{A}{(i), (iv)}
\opt{B}{(i), (iii)}
\opt{C}{(ii), (iii)}
\opt{D}{(ii), (iv)}
\end{qoptions}
\qn{3} Match the molecules under list (1) with shapes under list (2) select the correct answers from the sets (A), (B),(C) and (D) -\par\smallskip
{\footnotesize\setlength{\tabcolsep}{2pt}\begin{tabular}{|l|l|l|l|}
\hline
\multicolumn{2}{|c|}{List (1) molecule} & \multicolumn{2}{c|}{List (2) shape}\\
\hline
(P) & $\mathrm{PCl_5}$ & (i) & Linear\\
(Q) & $\mathrm{XeF_2}$ & (ii) & Triangular planar\\
(R) & $\mathrm{SO_3}$ & (iii) & Trigonal bipyramidal\\
(S) & $\mathrm{PH_3}$ & (iv) & Trigonal pyramidal\\
 & & (v) & Tetrahedral\\
\hline
\end{tabular}}
\begin{qoptions}
\opt{A}{P - (iv), Q - (i), R - (ii), S - (iii)}
\opt{B}{P - (iii), Q - (i), R - (ii), S - (v)}
\opt{C}{P - (iii), Q - (i), R - (v), S - (iv)}
\opt{D}{P - (iii), Q - (i), R - (ii), S - (iv)}
\end{qoptions}
\qn{4} \par\smallskip
{\footnotesize\setlength{\tabcolsep}{2pt}\begin{tabular}{|l|l|l|l|}
\hline
\multicolumn{2}{|c|}{Column-I (Complex ion)} & \multicolumn{2}{c|}{\shortstack{Column-II (Hybridisation,\\No. of unpaired electrons)}}\\
\hline
(P) & $\mathrm{[Cr(H_2O)_6]^{3+}}$ & (1) & $\mathrm{sp^3}$, 3\\
(Q) & $\mathrm{[CoCl_4]^{2-}}$ & (2) & $\mathrm{sp^3d^2}$, 5\\
(R) & $\mathrm{[Ni(NH_3)_6]^{2+}}$ & (3) & $\mathrm{d^2sp^3}$, 3\\
(S) & $\mathrm{[MnF_6]^{4-}}$ & (4) & $\mathrm{dsp^3}$, 3\\
 & & (5) & $\mathrm{sp^3d^2}$, 2\\
\hline
\end{tabular}}
\begin{qoptions}
\opt{A}{P$\to$3, Q$\to$1, R$\to$5, S$\to$2}
\opt{B}{P$\to$3, Q$\to$5, R$\to$1, S$\to$2}
\opt{C}{P$\to$1, Q$\to$3, R$\to$5, S$\to$2}
\opt{D}{P$\to$3, Q$\to$1, R$\to$2, S$\to$5}
\end{qoptions}
\qn{5} \par\smallskip
{\footnotesize\setlength{\tabcolsep}{2pt}\begin{tabular}{|l|l|l|l|}
\hline
\multicolumn{2}{|c|}{List-I (Complex)} & \multicolumn{2}{c|}{List-II (Property)}\\
\hline
(P) & $\mathrm{[Ni(CO)_4]}$ & (1) & Paramagnetic\\
(Q) & $\mathrm{[Cu(NH_3)_4]SO_4}$ & (2) & Diamagnetic\\
(R) & $\mathrm{K_3[Co(CN)_6]}$ & (3) & Square planar\\
(S) & $\mathrm{[Mn(H_2O)_6]Cl_2}$ & (4) & High spin\\
 & & (5) & Low spin\\
\hline
\end{tabular}}\par\smallskip
Select the most appropriate match :-
\begin{qoptions}
\opt{A}{P$\to$3, Q$\to$1, R$\to$4, S$\to$4}
\opt{B}{P$\to$1, Q$\to$2, R$\to$5, S$\to$5}
\opt{C}{P$\to$2, Q$\to$1, R$\to$5, S$\to$4}
\opt{D}{P$\to$2, Q$\to$3, R$\to$1, S$\to$5}
\end{qoptions}
\qn{6} Select the correct statement(s) about $\mathrm{[Fe(CN)_6]^{3-}}$ :-
\begin{qoptions}
\opt{A}{Total seven linkage isomers possible}
\opt{B}{Optical isomerism not possible in any possible linkage isomer.}
\opt{C}{Total possible isomers = 10}
\opt{D}{All of above are correct.}
\end{qoptions}
\qn{7} which of the following overlaps is incorrect (assuming Z-axis is internuclear axis)?\\
(i) $2p_y+2p_y\to\pi$-Bond formation\\
(ii) $2p_x+2p_x\to\sigma$-Bond formation\\
(iii) $3d_{xy}+3d_{xy}\to\pi$-Bond formation\\
(iv) $2s+2p_y\to\pi$-Bond formation\\
(v) $3d_{xy}+3d_{xy}\to\delta$-Bond formation\\
(vi) $2p_z+2p_z\to\sigma$-Bond formation
\begin{qoptions}
\opt{A}{(i), (ii), (iii)}
\opt{B}{(iii), (vi)}
\opt{C}{(ii), (v)}
\opt{D}{(ii), (iii), (iv)}
\end{qoptions}
\qn{8} A piston can freely move inside a horizontal cylinder closed from both ends. Initially, the piston separates the inside space of the cylinder into two equal parts each of volume $V_0$, in which an ideal gas is contained under the same pressure $P_0$ and at the same temperature. What work has to be performed in order to increase isothermally the volume of one part of gas $\eta$ times compared to that of the other by slowly moving the piston ?
\begin{qoptions}
\opt{A}{$P_0V_0\ln\eta$}
\opt{B}{$P_0V_0\ln\dfrac{(\eta+1)^2}{4\eta}$}
\opt{C}{$P_0V_0\ln\dfrac{(\eta-1)^2}{4\eta}$}
\opt{D}{$2P_0V_0\ln\eta$}
\end{qoptions}
\qn{9} At 35$^\circ$C, the vapour pressure of $\mathrm{CS_2}$ is 512 mm Hg and that of acetone $\mathrm{CH_3COCH_3}$, is 344 mm Hg. A solution of $\mathrm{CS_2}$ and acetone, in which the mole fraction of $\mathrm{CS_2}$ is 0.25 has a total vapour pressure of 600 mm Hg. Which the following statement about solution of acetone and $\mathrm{CS_2}$ is true ?
\begin{qoptions}
\opt{A}{A mixture of 100 ml of acetone and 100 ml of $\mathrm{CS_2}$ has a total volume of 200 ml.}
\opt{B}{When acetone and $\mathrm{CS_2}$ are mixed at 35$^\circ$C, heat is released to surroundings.}
\opt{C}{When acetone and $\mathrm{CS_2}$ are mixed at 35$^\circ$C, heat is absorbed from surroundings.}
\opt{D}{Raoult's law is obeyed by both $\mathrm{CS_2}$ and acetone for the solution in which the mole fraction of $\mathrm{CS_2}$ is 0.25.}
\end{qoptions}
\qn{10} For the given reaction :\\
$\mathrm{2A(s)+B(g)\rightleftharpoons C(g)+2D(g)+E(s)}$, in cylinder piston system the extent of reaction of $B$ was found to be 20\% at 300 K and 24\% at 500 K. The rate of backward reaction :
\begin{qoptions}
\opt{A}{Increases with increase in pressure and temperature.}
\opt{B}{Increases with increase in pressure and decrease in temperature.}
\opt{C}{Depends on temperature only and decreases with increases in temperature.}
\opt{D}{Depends on temperature only and increases with increase in temperature.}
\end{qoptions}
\qn{11} The composition of a sample of cuprous sulphide is found to be $\mathrm{Cu_{1.92}O_{1.00}}$, due to metal deficient defect. The molar ratio of $\mathrm{Cu^{2+}}$ and $\mathrm{Cu^+}$ ions in the crystal is :
\begin{qoptions}
\opt{A}{4 : 5}
\opt{B}{1 : 12.5}
\opt{C}{1 : 23}
\opt{D}{1 : 24}
\end{qoptions}
\qn{12} Methyl acetate was hydrolyzed in HCl aqueous solution at 298K, samples of reaction mixture were withdrawn at different time intervals and titrated with 0.12N NaOH sample whose volumes were obtained as :\par\smallskip
{\footnotesize\setlength{\tabcolsep}{4pt}\begin{tabular}{|l|c|c|c|}
\hline
Time(mins) & 0 & 25 & $\infty$\\
\hline
Volume of NaOH(mL) & 18 & 30 & 42\\
\hline
\end{tabular}}\par\smallskip
Half- life for the reaction below is :\\
\par\noindent\hfill{\footnotesize$\mathrm{CH_3COOCH_3(aq.)+H_2O(\ell)}\xrightarrow{\mathrm{H^\oplus}}\mathrm{CH_3COOH(aq.)+CH_3OH(aq.)}$}\hfill\null\par
\begin{qoptions}
\opt{A}{12.5 minutes}
\opt{B}{25 minutes}
\opt{C}{37.5 minutes}
\opt{D}{50 minutes}
\end{qoptions}
\qn{13} Certain amount of oxygen gas was passed over heated carbon, causing 80\% of oxygen added to be converted into CO gas, as per reaction\\
$\mathrm{2C(s)+O_2(g)\to2CO(g)}$\\
If the final gaseous mixture is allowed to effuse then ratio of rate of effusion of CO to $\mathrm{O_2}$ will be
\begin{qoptions}
\opt{A}{$\sqrt{\dfrac87}$}
\opt{B}{$\sqrt{\dfrac{8^3}7}$}
\opt{C}{$\sqrt{\dfrac{4^2\times8}7}$}
\opt{D}{$\sqrt{\dfrac{2^2\times8}7}$}
\end{qoptions}
\qn{14} Mark the correct option assuming monobromination in all reactions (include stereo isomers wherever possible)\par\smallskip
{\footnotesize\setlength{\tabcolsep}{2pt}\begin{tabular}{|l|l|l|p{3.3cm}|}
\hline
\multicolumn{2}{|c|}{List-I} & \multicolumn{2}{c|}{List-II}\\
\hline
(P) & \qfigi{p1_chem_s1_q14_1} & (1) & Even number of optically active products\\
\hline
(Q) & \qfigi{p1_chem_s1_q14_2} & (2) & Odd number of total products obtained after fractional distillation\\
\hline
(R) & \qfigi{p1_chem_s1_q14_3} & (3) & Even number of total products\\
\hline
(S) & \qfigi{p1_chem_s1_q14_4} & (4) & Even number of optically inactive products\\
\hline
\end{tabular}}
\begin{qoptions}
\opt{A}{P$\to$1,2,4 ; Q$\to$2,3,4 ; R$\to$1,3,4 ; S$\to$1,3,4}
\opt{B}{P$\to$1,3,4 ; Q$\to$1,2,3 ; R$\to$2,3,4 ; S$\to$2,3,4}
\opt{C}{P$\to$1,3,4 ; Q$\to$2,3,4 ; R$\to$2,3,4 ; S$\to$1,3,4}
\opt{D}{P$\to$2,3,4 ; Q$\to$1,2,4 ; R$\to$1,3,4 ; S$\to$2,3,4}
\end{qoptions}
\qn{15} \par
\qfig{p1_chem_s1_q15_1}
Identify relationship between produts ?
\begin{qoptions}
\opt{A}{Diastereo Isomers}
\opt{B}{Enantiomers}
\opt{C}{Functional Isomers}
\opt{D}{Geometrical Isomers}
\end{qoptions}
\qn{16} The correct sequence of reagents for the following conversion is
\qfig{p1_chem_s1_q16_1}
\begin{qoptions}
\opt{A}{\footnotesize$\xrightarrow[(1)]{\mathrm{Hg(CH_3COO)_2/H_2O}}\xrightarrow[(2)]{\mathrm{NaBH_4}}\xrightarrow[(3)]{\mathrm{Al_2O_3,\Delta}}\xrightarrow[(4)]{\mathrm{dil.H_2SO_4}}$}
\opt{B}{\footnotesize$\xrightarrow[(1)]{\mathrm{B_2H_6/THF}}\xrightarrow[(2)]{\mathrm{H_2O_2/\overset{\ominus}{O}H}}\xrightarrow[(3)]{\mathrm{dil.H_2SO_4}}\xrightarrow[(4)]{\mathrm{conc.H_2SO_4}}$}
\opt{C}{\footnotesize$\xrightarrow[(1)]{\mathrm{Dil.H_2SO_4}}\xrightarrow[(2)]{\mathrm{Al_2O_3,\Delta}}\xrightarrow[(3)]{\mathrm{B_2H_6/THF}}\xrightarrow[(4)]{\mathrm{H_2O_2/\overset{\ominus}{O}H}}$}
\opt{D}{\footnotesize$\xrightarrow[(1)]{\mathrm{Dil.H_2SO_4}}\xrightarrow[(2)]{\mathrm{conc.H_2SO_4/\Delta}}\xrightarrow[(3)]{\mathrm{Hg(CH_3COO)_2/H_2O}}\xrightarrow[(4)]{\mathrm{NaBH_4}}$}
\end{qoptions}
\qn{17} The structure of the starting material $X$ is :
\qfig{p1_chem_s1_q17_1}
\begin{qoptions}
\opt{A}{\qfigi{p1_chem_s1_q17_A}}
\opt{B}{\qfigi{p1_chem_s1_q17_B}}
\opt{C}{\qfigi{p1_chem_s1_q17_C}}
\opt{D}{\qfigi{p1_chem_s1_q17_D}}
\end{qoptions}
\qn{18} The end product of following reaction is :
\qfig{p1_chem_s1_q18_1}
\begin{qoptions}
\opt{A}{$\alpha$-diketone}
\opt{B}{$\alpha$-Hydroxy acid}
\opt{C}{1, 2-Diol}
\opt{D}{$\alpha$-Hydroxy ketone}
\end{qoptions}
\qn{19} Identify the reactions which is not correctly matched with their major product?
\begin{qoptions}
\opt{A}{\qfigi{p1_chem_s1_q19_A}$\xrightarrow[\mathrm{Pd/BaSO_4}]{\mathrm{H_2}}$\qfigi{p1_chem_s1_q19_3}}
\opt{B}{\qfigi{p1_chem_s1_q19_B}}
\opt{C}{\qfigi{p1_chem_s1_q19_C}$\xrightarrow{\mathrm{NaBH_4}}$\qfigi{p1_chem_s1_q19_6}}
\opt{D}{\qfigi{p1_chem_s1_q19_D}}
\end{qoptions}
\qn{20} \qfig{p1_chem_s1_q20_1}
\qfig{p1_chem_s1_q20_2}
Identify $P_3$
\begin{qoptions}
\opt{A}{\qfigi{p1_chem_s1_q20_A}}
\opt{B}{\qfigi{p1_chem_s1_q20_B}}
\opt{C}{\qfigi{p1_chem_s1_q20_C}}
\opt{D}{\qfigi{p1_chem_s1_q20_D}}
\end{qoptions}
\columnbreak
\sect{SECTION-II}
\qn{1} How many species have at least 1 lone pairs on its central atom.\\
$\mathrm{PCl_4^+}$, $\mathrm{BrF_5}$, $\mathrm{TeF_5^-}$, $\mathrm{I(CN)_2^-}$, $\mathrm{ClF_3}$, $\mathrm{CH_3^+}$, $\mathrm{CO_2}$
\qn{2} Out of the following. How many have correct IUPAC naming :\\
(1) $\mathrm{[Ni(CN)_4]^{2-}}$ - Tetracyanonickel(II) ion\\
(2) $\mathrm{[Pt(py)_4][PtCl_4]}$ - Tetrapyridineplatinum(II) tetrachorideplatinate(II)\\
(3) $\mathrm{[Ni(dmg)_2]}$ - Bis(dimethylglyoximato)nickel(II)\\
(4) $\mathrm{K_3[Fe(CN)_5NO]}$ - Potassium pentacyanonitrosylferrate(II)\\
(5) $\mathrm{[Fe(CO)_5]}$ - Pentacyanocaronylferrate(0)\\
(6) $\mathrm{K_2[HgI_4]}$ - Potassiumtetraiodidomercurate(II)\\
(7) $\mathrm{[Pt(NH_3)_4][CuCl_4]}$ - Tetraammineplatinum(IV) tetrachloridocuprate(II)\\
(8) $\mathrm{[Cu(gly)_2]}$ - Diglycinatecopper(II)\\
(9) $\mathrm{K_4[Fe(CN)_6]}$ - Potassium hexacyanidoferrate(II)\\
(10) $\mathrm{[Pt(NH_3)_6]Cl_4}$ - Hexaammineplatinum(IV) chloride.
\qn{3} Number of reactions giving either decarboxylation or dehydration or both on heating :
\qfig{p1_chem_s2_q03_1}
\qn{4} Vapour pressure of solution containing 6 g of a non-volatile solute in 180 g water is 20 torr. If 1 mole of water is further added, vapour pressure increases by 0.02 torr. The molar mass (in g/mol) of the non-volatile solute is.
\qn{5} Methane when solidified crystallizes in a structure having cubic unit cell with edge length of 0.6 nm. The density of solid $\mathrm{CH_4}$ is only slightly greater than the density of liquid $\mathrm{CH_4}$. If the density of liquid $\mathrm{CH_4}$ is 0.5 g/cm$^3$, the number of $\mathrm{CH_4}$ molecules per unit cell in solid crystal of methane is :
\qn{6} The decomposition of $\mathrm{A_2B_2}$ follows the rate equation $r=K[\mathrm{A_2B_2}]^2$.\\
$\mathrm{A_2B_2\xrightarrow{K}A_2+B_2}$. The kinetic data of the reaction is given below.\par\smallskip
{\footnotesize\setlength{\tabcolsep}{3pt}\begin{tabular}{|l|c|}
\hline
Initial conc. of $\mathrm{A_2B_2}$ (mole/lit) & 7\\
\hline
First half life of $\mathrm{A_2B_2}$ in sec. & 1440\\
\hline
\end{tabular}}\par\smallskip
If the initial concentration of $\mathrm{A_2B_2}$ were 56 m mol/lit, Calculate the time (in sec) taken for the conversion of $\mathrm{A_2B_2}$ to reach 60\% of initial value.\\
(Report your answer as $\dfrac t{1000}$ )
\qn{7} A cube shaped crystal of an alkali metal, 2 mm on an edge, was vaporised in a 800 mL flask. The pressure of the resulting vapour was 7.6 mm at 727$^\circ$ C. The structure of the solid metal is known to be body centred cubic. If volume of each unit cell of this crystal is [$P\times10^{-27}$ m$^3$] then find value of (3P) ? [Given : $R=0.08$ L-atm mol$^{-1}$K$^{-1}$, $N_A=6\times10^{23}$]
\qn{8} \qfig{p1_chem_s2_q08_1}

How many moles of Grignard reagent are consumed in above 1 mole of given compound
\qn{9} In how many of the following cases ratio of sigma bond and pi bond is 1 : 1\par\smallskip
{\footnotesize\setlength{\tabcolsep}{2pt}\begin{tabular}{|l|l|l|l|}
\hline
(i) $\mathrm{C_2(CN)_4}$ & (ii) $\mathrm{(CN)_2}$ & (iii) $\mathrm{SO_3}$ & (iv) $\mathrm{C_3O_2}$\\
\hline
(v) $\mathrm{CO_2}$ & (vi) $\mathrm{S_3O_9}$ & (vii) $\mathrm{P_4O_6}$ & (viii) $\mathrm{P_4O_{10}}$\\
\hline
\end{tabular}}
\qn{10} \qfig{p1_chem_s2_q10_1}
molecular mass of $A$ is
\partsec{3}{MATHEMATICS}
\sect{SECTION-I}
\qn{1} If $f(x)=\displaystyle\lim_{h\to0}\dfrac{e^{\sin^{-1}(x+h)}-e^{\sin^{-1}x}}h$, then equation of tangent at $x=0$ to the curve is given by -
\begin{qoptions}
\opt{A}{$y=2x+1$}
\opt{B}{$y=x+3$}
\opt{C}{$2x+y=1$}
\opt{D}{$y=x+1$}
\end{qoptions}
\qn{2} If $y^2=2ax$, $a>0$ and $xy=4\sqrt2$ intersect orthogonally, then `$a$' equals to
\begin{qoptions}
\opt{A}{$\dfrac12$}
\opt{B}{2}
\opt{C}{$\dfrac32$}
\opt{D}{4}
\end{qoptions}
\qn{3} Let $f(x)=xe^{-x(x+1)}$ is increasing in $\left(-a,\dfrac ab\right)$, where $a\in N$, then value of $(a+b)$ is
\begin{qoptions}
\opt{A}{3}
\opt{B}{1}
\opt{C}{$\dfrac12$}
\opt{D}{2}
\end{qoptions}
\qn{4} Range of $f(x)=k\sin x+\cos x$ is $[-5,5]$, then total number of integers between the two possible value of `$k$' is -
\begin{qoptions}
\opt{A}{11}
\opt{B}{10}
\opt{C}{9}
\opt{D}{5}
\end{qoptions}
\qn{5} If $\sin A+\cos B=a$ and $\sin B+\cos A=b$, then $\sin(A+B)$ is equal to -
\begin{qoptions}
\opt{A}{$\dfrac{a^2+b^2}2$}
\opt{B}{$\dfrac{a^2-b^2+2}2$}
\opt{C}{$\dfrac{a^2+b^2-2}2$}
\opt{D}{$\dfrac{a^2+b^2+2}2$}
\end{qoptions}
\qn{6} The value of $\cos\dfrac\pi5\cos\dfrac{2\pi}5\cos\dfrac{4\pi}5\cos\dfrac{8\pi}5$ is equal to -
\begin{qoptions}
\opt{A}{$\dfrac1{16}$}
\opt{B}{$\dfrac14$}
\opt{C}{$-\dfrac18$}
\opt{D}{$-\dfrac1{16}$}
\end{qoptions}
\qn{7} Number of solutions of the equation $\tan\theta+\tan2\theta+\tan3\theta=\tan\theta\tan2\theta\tan3\theta$ in $(0,\pi)$ is -
\begin{qoptions}
\opt{A}{5}
\opt{B}{3}
\opt{C}{2}
\opt{D}{zero}
\end{qoptions}
\qn{8} The number of solutions of $2^{\sin|x|}=4^{|\cos x|}$ in $[-\pi,\pi]$ is -
\begin{qoptions}
\opt{A}{0}
\opt{B}{2}
\opt{C}{4}
\opt{D}{6}
\end{qoptions}
\qn{9} The general solution of the equation $\sin x-\cos x=\min\left\{\dfrac1{\sqrt2},x^2-6x+10\right\}$ where $n\in I$, is -
\begin{qoptions}
\opt{A}{$\dfrac\pi4+n\pi+(-1)^n\dfrac\pi6$}
\opt{B}{$n\pi+(-1)^n\dfrac\pi6-\dfrac\pi4$}
\opt{C}{$2n\pi$}
\opt{D}{$2n\pi+\dfrac\pi2$}
\end{qoptions}
\qn{10} As shown in the figure, the area of the trapezium $ABCD$ is $\dfrac{81}2$, where $a=2b$, then $c$ is equal to -
\qfig{p1_math_s1_q10_1}
\begin{qoptions}
\opt{A}{$\dfrac92$}
\opt{B}{$3\sqrt5$}
\opt{C}{$\sqrt{30}$}
\opt{D}{none of these}
\end{qoptions}
\qn{11} A value of $c$ for which the Rolle's theorem holds for the function $f(x)=2x^3+x^2-4x-2$ on the interval $\left[-\sqrt2,\sqrt2\right]$ is
\begin{qoptions}
\opt{A}{$\dfrac35$}
\opt{B}{$\dfrac45$}
\opt{C}{$\dfrac23$}
\opt{D}{$-\dfrac23$}
\end{qoptions}
\qn{12} A spherical balloon is filled with $7500\pi$ m$^3$ of helium gas. If a leakage in the balloon starts at the rate of 100$\pi$ m$^3$/min., then the rate (in m/min) at which the radius of balloon will be decreasing after 30 minutes
\begin{qoptions}
\opt{A}{$\dfrac49$}
\opt{B}{$\dfrac29$}
\opt{C}{$\dfrac19$}
\opt{D}{$\dfrac13$}
\end{qoptions}
\qn{13} Which of the following inequality is incorrect ?
\begin{qoptions}
\opt{A}{$\ell n\left(1+\dfrac1x\right)>\dfrac2{2x+1}$ for $x>0$}
\opt{B}{$\sin x<x<\tan x$ for $x\in\left(0,\dfrac\pi2\right)$}
\opt{C}{$2x\sec x+x>3\tan x$ for $x\in\left(0,\dfrac\pi2\right)$}
\opt{D}{$\sin x<\dfrac{2x}\pi$ for $x\in\left(0,\dfrac\pi2\right)$}
\end{qoptions}
\qn{14} The function $f(x)=\begin{cases}\dfrac{|x|-x\left(3^{1/x}+1\right)}{3^{1/x}-1}&x\ne0\\0&x=0\end{cases}$ is
\begin{qoptions}
\opt{A}{discontinuous at $x=0$}
\opt{B}{continuous at $x=0$ but not differentiable at $x=0$}
\opt{C}{both continuous and differentiable at $x=0$}
\opt{D}{differentiable but not continuous at $x=0$}
\end{qoptions}
\qn{15} The value of $\displaystyle\lim_{\theta\to0^+}\left(\theta^2\sqrt\theta\cdot\sin\left(\dfrac1\theta\right)+4\sin\left(\theta\sqrt\theta\right)\right)^\theta$ is
\begin{qoptions}
\opt{A}{$-1$}
\opt{B}{1}
\opt{C}{2}
\opt{D}{$e$}
\end{qoptions}
\qn{16} $f(x)=x^5+e^{x/3}$, $g(x)=f^{-1}(x)$ then $g'(1)=$
\begin{qoptions}
\opt{A}{2}
\opt{B}{1}
\opt{C}{3}
\opt{D}{None of these}
\end{qoptions}
\qn{17} $y=e^{2\sqrt x}+e^{-2\sqrt x}$, then $\left(x\dfrac{d^2y}{dx^2}+\dfrac12\dfrac{dy}{dx}\right)+2$ equals to
\begin{qoptions}
\opt{A}{$y+2$}
\opt{B}{$xy+2$}
\opt{C}{$\dfrac{y+8}4$}
\opt{D}{$\sqrt{xy}$}
\end{qoptions}
\qn{18} Let\par\noindent\hfill{\small$f(x)=\begin{cases}x^2e^{2(x-1)}&;\ 0\le x\le1\\a\,\mathrm{sgn}(x+1)\cos(2x-2)+bx^2&,\ 1<x\le2\end{cases}$}\hfill\null\par
If $f(x)$ is differentiable at $x=1$, then (where $\mathrm{sgn}(x)$ denotes signum function of $x$)
\begin{qoptions}
\opt{A}{$a=-1$, $b=2$}
\opt{B}{$a=-3$, $b=-2$}
\opt{C}{$a=-3$, $b=4$}
\opt{D}{$a=-5$, $b=-4$}
\end{qoptions}
\qn{19} Which of the following function is discontinuous at $x=0$, if $f(0)=1$
\begin{qoptions}
\opt{A}{$f(x)=\dfrac{\sin\left(\sin^2x\right)}{\left(e^{|\sin x|}-1\right)\sqrt{x^2}}$; $(x\ne0)$}
\opt{B}{$f(x)=(1+x)(\mathrm{sgn}(x))+x^2$ ; $(x\ne0)$ (where $\mathrm{sgn}(x)$ is signum function)}
\opt{C}{$f(x)=\left(\ell n\left(1+\tan^2x\right)\right)\left(\mathrm{cosec}\left(\left(e^x-1\right)^2\right)\right)$; $(x\ne0)$}
\opt{D}{$f(x)=\left(1+\sin^2x\right)^{\mathrm{cosec}x}$; $(x\ne0)$}
\end{qoptions}
\qn{20} $\displaystyle\lim_{x\to-\infty}\left(\sqrt{x^2-x+1}-ax+2b\right)=0$, then $a=$ \underline{\hspace{1.2cm}}.
\begin{qoptions}
\opt{A}{$a=0$}
\opt{B}{$a=\pm2$}
\opt{C}{$a=-1$}
\opt{D}{$a=1$}
\end{qoptions}
\sect{SECTION-II}
\qn{1} If $f(x)=10x-2x^2$, then\par\noindent\hfill{\small$\displaystyle\lim_{h\to0}\dfrac{f(3\cosh+4\sinh-2)-8}{f\left(3e^h-5\,\mathrm{sech}\,h+4\right)-12}$}\hfill\null\par
is equal to
\qn{2} If $f(x)=\begin{cases}\dfrac{e^{ax}+\ell n(1+4x)+b}{x^3}&x\ne0\\L&x=0\end{cases}$ is continuous at $x=0$, then $12abL$ is equal to
\qn{3} Let $f(x)=(x-1)^{100}(x-2)^{2(99)}(x-3)^{3(98)}\ldots(x-100)^{100}$. If $P=\dfrac{f'(101)}{f(101)}$, then $\dfrac P{50}$ equals
\qn{4} Number of integers in the range of function $f(x)=\dfrac{100}{3\left(x^6+12x^4+5x^2+1\right)}$ is
\qn{5} If the solution set of inequality $\left(\cos x+\dfrac{\sqrt3}2\right)\left(\cos x-\dfrac12\right)\le0$ in $[0,\ 2\pi]$ is $\left[\dfrac{\alpha\pi}6,\dfrac{\beta\pi}6\right]\cup\left[\dfrac{\gamma\pi}6,\dfrac{\delta\pi}6\right]$, $\alpha,\beta,\gamma,\delta\in I^+$, then the value of $|\beta-\alpha+\delta-\gamma|$ is
\qn{6} If $x\in\left(\pi,\dfrac{3\pi}2\right)$, then value of $4\cos^2\left(\dfrac\pi4-\dfrac x2\right)+\sqrt{4\sin^4x+\sin^22x}$ is equal to
\qn{7} If $A+B+C=\pi$, and $\dfrac{\Sigma\sin2A}{\Sigma\sin A}=\lambda\prod\sin\dfrac A2$, then value of $\lambda$ is
\qn{8} If $f(x)+f\left(1-\dfrac1x\right)=1+x\ \forall\,x\in R-\{0,1\}$, then value of $16f(2)$ is
\qn{9} For $0<x<1$, if the solution of equation $3\sin^{-1}\left(\dfrac{2x}{1+x^2}\right)+2\tan^{-1}\left(\dfrac{2x}{1-x^2}\right)-4\cos^{-1}\left(\dfrac{1-x^2}{1+x^2}\right)=\dfrac\pi4$ is `$K$' then value of `$100K$' is (use $\sqrt2=1.41$)
\qn{10} If $\displaystyle\lim_{x\to2}\dfrac{1-\cos\left(x^2-6x+8\right)}{(x-2)^2}=p$, then value of $\left(\dfrac6p\right)$ is equal to
\raggedcolumns
\papersec{2}
\partsec{1}{PHYSICS}
\sect{SECTION-I(i)}
\qn{1} The figure shows the variation of electrostatic potential $V$ in volt with the distance of position of point along $x$ axis from origin due to continuous volume charge distribution in the region $x=-1$m to $x=+1$m, the graph is parabolic ($V=15-5x^2$) and rest portion of graph is straight line. Choose the CORRECT option(s)\\
($\varepsilon_0\to$ permittivity of free space $=8.85\times10^{-12}\ \mathrm{N^{-1}m^{-2}C^2}$). The direction of $\vec E$ along positive $x$-axis is considered as positive.
\qfig{p2_phy_s1_q01_1}
\begin{qoptions}
\opt{A}{This graph of potential may be due to a thick sheet of infinite dimension ($-1\mathrm m\le x\le1$, $-\infty<y<\infty$ and $-\infty<z<\infty$) with constant volume charge density $1.77\times10^{-19}$ C/m$^2$.}
\opt{B}{This graph of potential may be due to a thick sheet of infinite dimension ($-1\mathrm m\le x\le1$, $-\infty<y<\infty$ and $-\infty<z<\infty$) with constant volume charge density $0.885\times10^{-10}$ C/m$^3$.}
\opt{C}{\qfigi{p2_phy_s1_q01_C}}
\opt{D}{\qfigi{p2_phy_s1_q01_D}}
\end{qoptions}
\qn{2} The given two fixed rings of radius $R$ lie in $XY$ plane as shown in the figure. Linear charge density of the rings $A$ and $B$ varies as per the relation $\lambda=\lambda_0\sin\theta$ and $\lambda=\lambda_0\cos\theta$ respectively where $\theta$ is measured from $+x$ axis. Distance between the centre of the two rings is $r\gg R$.
\qfig{p2_phy_s1_q02_1}
\begin{qoptions}
\opt{A}{Force of interaction between the two ring is $\dfrac34\times\dfrac{\lambda_0^2\pi}{\in_0}\left(\dfrac Rr\right)^4$.}
\opt{B}{Force of interaction between the two ring is $\dfrac32\times\dfrac{\lambda_0^2\pi}{\in_0}\left(\dfrac Rr\right)^4$.}
\opt{C}{Potential energy of interaction between the two ring is $\dfrac{\pi\lambda_0^2R^4}{4\varepsilon_0r^3}$.}
\opt{D}{Potential energy of interaction between the two ring is zero.}
\end{qoptions}
\qn{3} A point charge $q$ is placed at $P(0,0,a)$. Then select correct alternatives.
\qfig{p2_phy_s1_q03_1}
\begin{qoptions}
\opt{A}{$\dfrac q{24\varepsilon_0}>\phi_{OAB}>\dfrac q{48\varepsilon_0}$}
\opt{B}{$\phi_{OAC}=\dfrac q{48\varepsilon_0}$}
\opt{C}{$\phi_{ABC}<\phi_{ABO}$}
\opt{D}{$\phi_{OAC}<\phi_{OBC}$}
\end{qoptions}
\qn{4} A satellite revolves around a spherical planet in a stable elliptic orbit. Its speed at perigee is 2 km/s and its speed at apogee is 1 km/s. Its farthest distance from the planet's centre is 6000 km. Then which of the following statements are correct.
\begin{qoptions}
\opt{A}{The eccentricity of its orbit is 1/2.}
\opt{B}{The line joining the planet's center to the satellite sweeps out an area of $1.8\times10^{11}$ m$^2$ in each minute.}
\opt{C}{The planet's radius must be less than 3000 km.}
\opt{D}{The gain in its potential energy when it travels from the perigee to apogee is 3/4 times its kinetic energy at the perigee.}
\end{qoptions}
\qn{5} Two satellites of same mass of a planet in circular orbits have periods of revolution 32 days and 256 days. If the radius of the orbit of the first is $R$, then the
\begin{qoptions}
\opt{A}{Radius of the orbit of the second is 8$R$}
\opt{B}{Radius of the orbit of the second is 4$R$}
\opt{C}{Total mechanical energy of the second is greater than that of the first.}
\opt{D}{Kinetic energy of the second is less than that of the first.}
\end{qoptions}
\qn{6} In the potentiometer arrangement shown, the driving cell $D$ has emf $\varepsilon$ and internal resistance $r$. The cell $C$, whose emf is to be measured, has emf $\varepsilon/2$ and internal resistance $2r$. The potentiometer wire is 100-cm long. If balance is obtained at the length $AJ=\ell$.
\qfig{p2_phy_s1_q06_1}
\begin{qoptions}
\opt{A}{$\ell=50$ cm}
\opt{B}{$\ell>50$ cm}
\opt{C}{Balance will be obtained only if resistance of $AB$ is $\ge r$}
\opt{D}{Balance cannot be obtained}
\end{qoptions}
\qn{7} In the circuit shown $AB$ is potentiometer wire of length 10 m and resistance 500 $\Omega$. $CD$ is a 1 m wire of meter bridge which is balanced at $CP=20$ cm. Given that when the potential difference across $R_1$ is applied across $XY$ the balancing length on potentiometer is 2 m and for potential difference across $R_2$, the corresponding length is 8 m. Then :-
\qfig{p2_phy_s1_q07_1}
\qfig{p2_phy_s1_q07_2}
\begin{qoptions}
\opt{A}{$\dfrac{R_1}{R_2}=\dfrac14$}
\opt{B}{Potential difference across $R_1$ is 1V}
\opt{C}{Potential difference across $R_2$ is 4V}
\opt{D}{EMF of the battery $E=5$V}
\end{qoptions}
\sect{SECTION-I(ii)}
\paraq{8}{10}
Three columns are given. Column--I shows conductor in different situation, whereas column--II and column-III represents magnitude of electric field and modulus of potential at point `$P$' respectively.
\end{multicols}
\noindent\begin{minipage}{\textwidth}
\footnotesize\setlength{\tabcolsep}{4pt}
\begin{tabular}{|l|p{6.6cm}|l|l|l|l|}
\hline
\multicolumn{2}{|c|}{Column-I} & \multicolumn{2}{c|}{Column-II} & \multicolumn{2}{c|}{Column-III}\\
\hline
(I) & Electric field \& potential at `$P$' due to charges on shell is $E$ \& $V$ respectively.\par\qfigi{p2_phy_s2_para8_10_1} & (i) & $E=\dfrac{5Kq}{8R^2}$ & (P) & $V=\dfrac{15Kq}{6R}$\\
\hline
(II) & Solid conductor of inner radius $R$ and outer radius $2R$. Net Electric field \& potential at `$P$' is $E$ \& $V$ respectively.\par\qfigi{p2_phy_s2_para8_10_2} & (ii) & $E=\dfrac{4Kq}{9R^2}$ & (Q) & $V=\dfrac{4Kq}{6R}$\\
\hline
(III) & Electric field \& potential at `$P$' due to charges on shell is $E$ \& $V$ respectively.\par\qfigi{p2_phy_s2_para8_10_3} & (iii) & $E=\dfrac{2Kq}{9R^2}$ & (R) & $V=\dfrac{17Kq}{6R}$\\
\hline
(IV) & $P$ point lies on the axis of co-axial rings having radius $R$ and $\sqrt3R$ \& charges $\sqrt2q$ \& $q$ respectively.\par\qfigi{p2_phy_s2_para8_10_4} & (iv) & $E=\dfrac{4Kq}{R^2}$ & (S) & $V=\dfrac{9Kq}{6R}$\\
\hline
\end{tabular}
\end{minipage}
\begin{multicols}{2}
\qn{8} Choose the correct option:
\begin{qoptions}
\opt{A}{(IV)-(i)-(Q)}
\opt{B}{(III)-(iv)-(P)}
\opt{C}{(III)-(iii)-(Q)}
\opt{D}{None}
\end{qoptions}
\qn{9} Choose the correct option:
\begin{qoptions}
\opt{A}{(II)-(i)-(Q)}
\opt{B}{(I)-(ii)-(R)}
\opt{C}{(I)-(iii)-(Q)}
\opt{D}{None}
\end{qoptions}
\qn{10} Choose the correct option:
\begin{qoptions}
\opt{A}{(II)-(iv)-(P)}
\opt{B}{(II)-(iii)-(P)}
\opt{C}{(II)-(iii)-(Q)}
\opt{D}{None}
\end{qoptions}
\paraq{11}{13}
A particle is taken to a distance $r(>R)$ from centre of the earth. $R$ is radius of the earth. It is given velocity $V$ which is perpendicular to $\vec r$. With the given values of $V$ in column I you have to match the values of total energy of particle in column II and the resultant path of particle in column III. Here `$G$' is the universal gravitational constant and `$M$' is the mass of the earth.\par\smallskip
{\footnotesize\setlength{\tabcolsep}{1pt}\begin{tabular}{|l|l|l|l|l|l|}
\hline
 & \shortstack{Column-I\\(Velocity)} & & \shortstack{Column-II\\(Total Energy)} & & \shortstack{Column-III\\(Path)}\\
\hline
(I) & $V=\sqrt{\dfrac{GM}r}$ & (i) & Negative & (P) & Elliptical\\
(II) & $V=\sqrt{\dfrac{2GM}r}$ & (ii) & Positive & (Q) & Parabolic\\
(III) & $V>\sqrt{\dfrac{2GM}r}$ & (iii) & Zero & (R) & Hyperbolic\\
(IV) & $\sqrt{\dfrac{GM}r}<V<\sqrt{\dfrac{2GM}r}$ & (iv) & Infinite & (S) & Circular\\
\hline
\end{tabular}}
\qn{11} Choose the correct option :
\begin{qoptions}
\opt{A}{(I)-(i)-(S)}
\opt{B}{(II)-(iv)-(R)}
\opt{C}{(III)-(ii)-(Q)}
\opt{D}{(IV)-(iii)-(P)}
\end{qoptions}
\qn{12} Choose the correct option :
\begin{qoptions}
\opt{A}{(I)-(i)-(R)}
\opt{B}{(II)-(iii)-(Q)}
\opt{C}{(III)-(ii)-(S)}
\opt{D}{(IV)-(iv)-(P)}
\end{qoptions}
\qn{13} Choose the correct option :
\begin{qoptions}
\opt{A}{(III)-(ii)-(R)}
\opt{B}{(II)-(i)-(S)}
\opt{C}{(I)-(iii)-(P)}
\opt{D}{(IV)-(iv)-(Q)}
\end{qoptions}
\sect{SECTION-III}
\qn{1} In the circuit shown in figure each battery is 5 V and has an internal resistance of 0.2 $\Omega$. What is the reading of ideal voltmeter $V$ (in Volt )?
\qfig{p2_phy_s3_q01_1}
\qn{2} The energy required to dislodge 25\% mass of a planet of mass $M$ and radius $R$, so that is can form a satellite orbiting at a radius of 1.5 $R$ from the centre of the planet, is given by $E=\dfrac{0.21}x\dfrac{GM^2}R$. What is the nearest integer value of $x$? Assume that both the bodies are of spherical shape. [$G$ -- universal gravitational constant. Take $4^{1/3}=1.6$ and $\left(\dfrac43\right)^{1/3}=1.1$]
\qn{3} A particle of mass $m$ is moving in circular motion of radius $r$ with constant speed $v$ in influence of two fixed particles of mass $M$. Plane of circular motion is perpendicular to line joining fixed particles as shown. Speed of particle is given as $v=\sqrt{\dfrac{GM}{nr}}$ find value of $n$ ($a=r\sqrt3$)
\qfig{p2_phy_s3_q03_1}
\qn{4} A ring of mass $m$ and radius $3R$ is rotating with constant angular speed $\omega$ around a planet of mass $M$ and radius $R$. Center of ring and planet concide with each other. Tension in the ring is given as $T=\dfrac{GMm}{3n\pi R^2}$. Find value of $n$ $\left(\omega=\sqrt{\dfrac{GM}{9R^3}}\right)$.
\qfig{p2_phy_s3_q04_1}
\qn{5} A solid sphere of radius `$R$' is uniformly charged with charge density $\rho$ in its volume. A spherical cavity of radius $\dfrac R2$ is made in the sphere as shown in the figure. The electric potential at the centre of the sphere is $\dfrac{n\rho R^2}{12\in_0}$. Find the value of $n$.
\qfig{p2_phy_s3_q05_1}
\partsec{2}{CHEMISTRY}
\sect{SECTION-I(i)}
\qn{1} Which of the following reaction will generate 3$^\circ$ alcohol?
\begin{qoptions}
\opt{A}{\qfigi{p2_chem_s1_q01_A}}
\opt{B}{\qfigi{p2_chem_s1_q01_B}}
\opt{C}{\qfigi{p2_chem_s1_q01_C}}
\opt{D}{\qfigi{p2_chem_s1_q01_D}}
\end{qoptions}
\qn{2} Correct option amongst the following is/are :
\begin{qoptions}
\opt{A}{\qfigi{p2_chem_s1_q02_A}}
\opt{B}{\qfigi{p2_chem_s1_q02_B}}
\opt{C}{\qfigi{p2_chem_s1_q02_C}}
\opt{D}{\qfigi{p2_chem_s1_q02_D}}
\end{qoptions}
\qn{3} Spinel is an important class of oxides consisting of two types of metal ions with the oxide ions arranged in CCP pattern. The normal spinel has one-eighth of the tetrahedral holes in the crystal structure occupied by one type of metal ions and one-half of the octahedral holes in the crystal structure occupied by another type of metal ions. One such type of a spinel is formed by $\mathrm{Zn^{2+}}$, $\mathrm{Al^{3+}}$ and $\mathrm{O^{2-}}$ ions. The simplest formula of this spinel can not be :
\begin{qoptions}
\opt{A}{$\mathrm{ZnAl_2O_4}$}
\opt{B}{$\mathrm{Zn_2AlO_4}$}
\opt{C}{$\mathrm{Zn_2Al_3O_4}$}
\opt{D}{$\mathrm{ZnAlO_2}$}
\end{qoptions}
\qn{4} Which of the following plot(s) is/are correct for a sample of ideal gas in given conditions ?
\begin{qoptions}
\opt{A}{\qfigi{p2_chem_s1_q04_A}}
\opt{B}{\qfigi{p2_chem_s1_q04_B}}
\opt{C}{\qfigi{p2_chem_s1_q04_C}}
\opt{D}{\qfigi{p2_chem_s1_q04_D}}
\end{qoptions}
\qn{5} Werner synthesized and resolved a polynuclear complex containing no carbon atom, Tris[tetraamine-$\mu$-dihydroxidocobalt(III)]cobalt(III) cation also called Hexol.\\
Select the correct statement(s) regarding given complex Hexol.
\qfig{p2_chem_s1_q05_1}
\begin{qoptions}
\opt{A}{Total 3 octahedral units in one complex ion.}
\opt{B}{Octahedral units joined by common edges between two units.}
\opt{C}{Shows optical isomerism.}
\opt{D}{Total 6 bridging hydroxide ligands.}
\end{qoptions}
\qn{6} Consider the given set of species and select correct statements.\\
$\mathrm{H_2}$, $\mathrm{Li_2}$, $\mathrm{B_2}$, $\mathrm{C_2}$, $\mathrm{N_2}$, $\mathrm{O_2}$, $\mathrm{F_2}$\\
Condition 1 $\Rightarrow$ Diamagnetic molecule\\
Condition 2 $\Rightarrow$ Last $\mathrm{e^-}$ present in gerade B.M.O.\\
Condition 3 $\Rightarrow$ On $\mathrm{e^-}$ addition stability decreases.
\begin{qoptions}
\opt{A}{Only three species follows all three conditions}
\opt{B}{Only two species follows only two conditions}
\opt{C}{Only one specie follows only one condition}
\opt{D}{Only one specie that do not follow any of above condition}
\end{qoptions}
\qn{7} Select correct option for given options.\\
(i) Glycinate (gly)\\
(ii) Acelylacetanato ($\mathrm{acac^-}$)\\
(iii) Dipyridine (dipy)\\
(iv) Butylenediamine (bn)\\
(v) Ethylenediamine (en)\\
(vi) 1,10-phenthroline (o-phen)\\
(vii) Acetate ($\mathrm{OAc^-}$)\\
(viii) EDTA$^{4-}$ (Ethylenediaminetetraacetate)\\
(ix) Diethylene triamine (dien)
\begin{qoptions}
\opt{A}{Six ligands form 5 member ring with atleast 2 nitrogen atoms and metal centre}
\opt{B}{One ligands form 5 member ring with two oxygen atom and metal centre}
\opt{C}{Seven ligands form 5 member ring with 2 carbon atoms and metal centre}
\opt{D}{Three ligands already contain 6 member ring.}
\end{qoptions}
\sect{SECTION-I(ii)}
\paraq{8}{10}
Initially `$a$' moles of $A$ react with 2 moles of $B$ in a 1 litre container according to given reactions for different values of `$a$' :
\end{multicols}
\noindent\begin{minipage}{\textwidth}
\footnotesize\setlength{\tabcolsep}{4pt}
\begin{tabular}{|l|p{5.2cm}|l|p{3.2cm}|l|p{4.6cm}|}
\hline
\multicolumn{2}{|c|}{Column-I} & \multicolumn{2}{c|}{Column-II} & \multicolumn{2}{c|}{Column-III}\\
\hline
(I) & $\mathrm{A(g)+2B(g)\to C(g)+D(g)}$\newline $r=(2\times10^{-3})[A]^{1/2}[B]^{1/2}$\newline $a=1$ & (i) & Has highest value of half life & (P) & May be an elementary reaction\\
\hline
(II) & $\mathrm{2A(g)+B(g)\to C(g)+2D(g)}$\newline $r=(2\times10^{-3})[A]^{1}[B]^{0}$\newline $a=4$ & (ii) & Has lowest value of half life & (Q) & Must be a complex reaction\\
\hline
(III) & $\mathrm{A(g)+B(g)\to 2C(g)+D(g)}$\newline $r=(2\times10^{-3})[A][B]$\newline $a=2$ & (iii) & Reaction is complete in definite time & (R) & After equal time intervals remaining concentrations of $B$ will be in geometric progression\\
\hline
(IV) & $\mathrm{A+3B\to 2C(g)+3D(g)}$\newline $r=(2\times10^{-3})[A]^{0}[B]^{2}$\newline $a=(2/3)$ & (iv) & Reaction is complete in infinite time & (S) & After equal time intervals remaining concentration of $B$ will be in harmonic progression\\
\hline
\end{tabular}
\end{minipage}
\begin{multicols}{2}
\qn{8} For given reactions only correct option will be :
\begin{qoptions}
\opt{A}{(I) (ii) (P)}
\opt{B}{(II) (iii) (R)}
\opt{C}{(III) (i) (P)}
\opt{D}{(IV) (ii) (P)}
\end{qoptions}
\qn{9} For given reactions only incorrect option will be :
\begin{qoptions}
\opt{A}{(I) (iii) (Q)}
\opt{B}{(II) (iv) (R)}
\opt{C}{(III) (iv) (P)}
\opt{D}{(IV) (ii) (Q)}
\end{qoptions}
\qn{10} For (I) correct option will be :
\begin{qoptions}
\opt{A}{(I) (iii) (Q)}
\opt{B}{(I) (iv) (R)}
\opt{C}{(I) (i) (S)}
\opt{D}{(I) (ii) (Q)}
\end{qoptions}
\paraq{11}{13}
\end{multicols}
\noindent\begin{minipage}{\textwidth}
\footnotesize\setlength{\tabcolsep}{3pt}\begin{tabular}{|l|l|l|l|l|l|}
\hline
\multicolumn{2}{|c|}{Column-I} & \multicolumn{2}{c|}{Column-II} & \multicolumn{2}{c|}{Column-III}\\
\hline
(a) & $\mathrm{[Pt(NH_3)_4Cl_2]Br_2}$ & (i) & Inner orbital complex & (p) & Total possible isomers = 3\\
(b) & $\mathrm{[Co(en)(NH_3)_2Cl_2]Cl}$ & (ii) & Low spin complex & (q) & Total possible isomers = 6\\
(c) & $\mathrm{K_2[Ni(CN)_4]}$ & (iii) & Outer orbital complex & (r) & Total possible isomers = 8\\
(d) & $\mathrm{K_4[Ni(SCN)_6]}$ & (iv) & High spin complex & (s) & Total possible isomers = 10\\
\hline
\end{tabular}
\end{minipage}
\begin{multicols}{2}
\qn{11} Select the CORRECT combination is :
\begin{qoptions}
\opt{A}{a-ii-p}
\opt{B}{b-iv-q}
\opt{C}{c-i-s}
\opt{D}{d-iii-s}
\end{qoptions}
\qn{12} Select INCORRECT combination is :
\begin{qoptions}
\opt{A}{a-ii-q}
\opt{B}{a-i-q}
\opt{C}{c-ii-q}
\opt{D}{c-i-r}
\end{qoptions}
\qn{13} Select CORRECT combination in which maximum number of unpaired electrons present in complex :
\begin{qoptions}
\opt{A}{b-i-p}
\opt{B}{c-iii-q}
\opt{C}{d-iii-s}
\opt{D}{d-iv-s}
\end{qoptions}
\sect{SECTION-III}
\qn{1} How many reagents can be used for the given conversion.
\qfig{p2_chem_s3_q01_1}
(1) Cu, 300$^\circ$C \quad (2) $\mathrm{KMnO_4/H^\oplus}$ \quad (3) PCC/$\mathrm{CH_2Cl_2}$ \quad (4) $\mathrm{K_2Cr_2O_7/H^\oplus}$\\
(5) $\mathrm{MnO_2}$ \quad (6) $\mathrm{Al(OCMe_3)_3}$/Acetone\\
(7) NBS \quad (8) (i) TsCl (ii) DMSO (iii) $\mathrm{NaHCO_3}$\\
(9) (i) DIBAL-H (ii) $\mathrm{H_2O}$
\qn{2} How many reaction represent the correct major product :
\qfig{p2_chem_s3_q02_1}
(ii) \qfigi{p2_chem_s3_q02_2}\par
\qfig{p2_chem_s3_q02_3}
\qfig{p2_chem_s3_q02_4}
\qfig{p2_chem_s3_q02_5}
\qfig{p2_chem_s3_q02_6}
\qfig{p2_chem_s3_q02_7}
\qfig{p2_chem_s3_q02_8}
\qfig{p2_chem_s3_q02_9}
\qn{3} A strong current of trivalent gaseous boron passed through a silicon crystal decreases the density of the crystal due to part replacement of silicon by boron and due to interstitial vacancies created by missing Si atoms. In one such experiment, one gram of silicon is taken and the boron atoms are found to be 1000 ppm, by mass, when the density of the Si crystal decreases by 12\%. The percentage of missing vacancies due to silicon, which are filled up by boron atoms is (Atomic masses : Si = 30, B = 11)\\
\textbf{[Give your answer as nearest integer]}
\qn{4} In the given reaction sequence, find number of possible 5 and 6 membered Lactone(s) having different boiling points :
\qfig{p2_chem_s3_q04_1}
\qn{5} \par\smallskip
{\setlength{\tabcolsep}{3pt}\begin{tabular}{ll}
(a) \qfigi{p2_chem_s3_q05_1} & \qfigi{p2_chem_s3_q05_2}\\
(c) \qfigi{p2_chem_s3_q05_3} & (d) \qfigi{p2_chem_s3_q05_6}\\
(e) \qfigi{p2_chem_s3_q05_4} & \qfigi{p2_chem_s3_q05_5}\\
(g) \qfigi{p2_chem_s3_q05_9} & \qfigi{p2_chem_s3_q05_7}\\
\qfigi{p2_chem_s3_q05_8} & \qfigi{p2_chem_s3_q05_10}\\
\end{tabular}}\par\smallskip
$X=$ Number of compounds giving (+) test with $\mathrm{AgNO_3/NH_4OH}$\\
$Y=$ Number of compounds giving yellow colour with $\mathrm{I_2+NaOH}$\\
Find $X-Y=$ ?
\partsec{3}{MATHEMATICS}
\sect{SECTION-I(i)}
\qn{1} If $e^{\sin\left(x^2+y^2\right)}=\tan\dfrac{y^2}4+\sin^{-1}x$, then $y'(0)$ can be -
\begin{qoptions}
\opt{A}{$\dfrac1{3\sqrt\pi}$}
\opt{B}{$-\dfrac1{3\sqrt\pi}$}
\opt{C}{$-\dfrac1{5\sqrt\pi}$}
\opt{D}{$-\dfrac1{3\sqrt{5\pi}}$}
\end{qoptions}
\qn{2} If $\cot\theta+\tan\theta=x$ and $\sec\theta-\cos\theta=y$ wherever defined the which of following is/are true ?
\begin{qoptions}
\opt{A}{$x\sin\theta\cos\theta=1$}
\opt{B}{$\sin^2\theta=y\cos\theta$}
\opt{C}{$\left(x^2y\right)^{2/3}+\left(xy^2\right)^{2/3}=1$}
\opt{D}{$\left(x^2y\right)^{2/3}-\left(xy^2\right)^{2/3}=1$}
\end{qoptions}
\qn{3} For the curve $\sin x+\sin y=1$, $\displaystyle\lim_{x\to0}\left(x^\alpha\right)\dfrac{d^2y}{dx^2}$ exists then $\alpha$ can be
\begin{qoptions}
\opt{A}{$\dfrac32$}
\opt{B}{2}
\opt{C}{$\dfrac52$}
\opt{D}{1}
\end{qoptions}
\qn{4} Let $x^{\frac23}+y^{\frac23}=a^{\frac23}$, ($a>0$) be the equation of curve then which of the following is/are correct ?
\begin{qoptions}
\opt{A}{Equation of tangent at $(x_1,y_1)$ on the curve is $\dfrac x{x_1^{1/3}}+\dfrac y{y_1^{1/3}}=a^{2/3}$}
\opt{B}{Length of portion of tangent intercepted between the coordinate axes is constant.}
\opt{C}{If tangent at any point on the curve meet the coordinate axes in $A$ and $B$, then locus of midpoint of $AB$ is $x^2+y^2=\dfrac{a^2}4$}
\opt{D}{Equation of normal at $(x_1,y_1)$ on the curve is $xx_1^{1/3}-yy_1^{1/3}=x_1^{4/3}+y_1^{4/3}$}
\end{qoptions}
\qn{5} If the curves $y=2(x-a)^2$ and $y=e^{2x}$ touches each other, then `$a$' is less than -
\begin{qoptions}
\opt{A}{$-1$}
\opt{B}{0}
\opt{C}{1}
\opt{D}{2}
\end{qoptions}
\qn{6} Let $S_n=\displaystyle\sum_{k=0}^{n}\tan^{-1}\dfrac{49}{49k^2+63k+57}$ then $\cot\left(\displaystyle\lim_{n\to\infty}S_n\right)$ is equal to
\begin{qoptions}
\opt{A}{0.24}
\opt{B}{0.14}
\opt{C}{24}
\opt{D}{14}
\end{qoptions}
\qn{7} Given that $f(x)$ is a function defined on $R$, satisfying $f(1)=1$ and for any $x\in R$, $f(x+5)\ge f(x)+5$ and $f(x+1)\le f(x)+1$. If $g(x)=f(x)+5-x$, then $g(2023)$ is greater than or equal to -
\begin{qoptions}
\opt{A}{3}
\opt{B}{4}
\opt{C}{5}
\opt{D}{6}
\end{qoptions}
\sect{SECTION-I(ii)}
\paraq{8}{10}
Columns I, II and III contain function, domain and range respectively.\par\smallskip
\end{multicols}
\noindent\begin{minipage}{\textwidth}
\footnotesize\setlength{\tabcolsep}{3pt}\begin{tabular}{|l|l|l|l|l|l|}
\hline
\multicolumn{2}{|c|}{Column-I} & \multicolumn{2}{c|}{Column-II} & \multicolumn{2}{c|}{Column-III}\\
\hline
(I) & $f(x)=\cos^{-1}\left(\dfrac{1-x^2}{1+x^2}\right)$ & (i) & $(-7,\infty)-\{9\}$ & (P) & $\left(\dfrac{-\pi}2,\dfrac\pi2\right)$\\
(II) & $g(x)=\tan^{-1}\left(\dfrac{3x-x^3}{1-3x^2}\right)$ & (ii) & $(-\infty,\infty)$ & (Q) & $[0,\pi)$\\
(III) & $h(x)=\pi\left(\dfrac{\sqrt{x+7}-4}{x-9}\right)$ & (iii) & $R-\left\{\dfrac{-1}{\sqrt3},\dfrac1{\sqrt3}\right\}$ & (R) & $[-\pi,\pi]$\\
(IV) & $k(x)=\dfrac\pi{\sqrt2}\left(\sin\sqrt{x^2}+\cos\sqrt{x^2}\right)$ & (iv) & $[0,\infty)$ & (S) & $\left(0,\dfrac\pi4\right]-\left\{\dfrac\pi8\right\}$\\
\hline
\end{tabular}
\end{minipage}
\begin{multicols}{2}
\qn{8} Which of the following options is the only CORRECT combination ?
\begin{qoptions}
\opt{A}{(I) (iv) (Q)}
\opt{B}{(III) (i) (P)}
\opt{C}{(IV) (ii) (R)}
\opt{D}{(II) (iii) (S)}
\end{qoptions}
\qn{9} Which of the following options is the only CORRECT combination ?
\begin{qoptions}
\opt{A}{(II) (ii) (P)}
\opt{B}{(I) (ii) (Q)}
\opt{C}{(III) (i) (R)}
\opt{D}{(IV) (iv) (R)}
\end{qoptions}
\qn{10} Which of the following options is the only CORRECT combination ?
\begin{qoptions}
\opt{A}{(IV) (ii) (Q)}
\opt{B}{(I) (iv) (Q)}
\opt{C}{(II) (iii) (R)}
\opt{D}{(III) (i) (S)}
\end{qoptions}
\paraq{11}{13}
\end{multicols}
\noindent\begin{minipage}{\textwidth}
\footnotesize\setlength{\tabcolsep}{3pt}\begin{tabular}{|l|p{6.5cm}|l|l|l|p{3.2cm}|}
\hline
\multicolumn{2}{|c|}{Column-I} & \multicolumn{2}{c|}{Column-II} & \multicolumn{2}{c|}{Column-III}\\
\hline
(I) & $f(x)=\mathrm{sgn}\left(\dfrac{|\sin x|}{\sin x}\right)$ (where $\mathrm{sgn}(.)$ is signum function) & (i) & even & (P) & maximum value of $f(x)=1$\\
\hline
(II) & $f(x)=\dfrac12\left(\{x\}+\left\{x+\dfrac73\right\}+\left\{x+\dfrac{17}3\right\}\right)$ (where $\{.\}$ represents fractional part function) & (ii) & odd & (Q) & value of $f(x)$ is always less than 1\\
\hline
(III) & $f(x)=\dfrac2\pi\left[\max\left\{\left|\dfrac\pi2-\cos^{-1}\cos x\right|,|\sin x|\right\}\right]$ & (iii) & periodic & (R) & minimum value of $f(x)=0$\\
\hline
(IV) & $f(x)=\dfrac1{|x-1|+|x|+|x+1|}$ & (iv) & bounded & (S) & minimum value of $|f(x)|=1$\\
\hline
\end{tabular}
\end{minipage}
\begin{multicols}{2}
\qn{11} Which of the following options is the only CORRECT combination ?
\begin{qoptions}
\opt{A}{(I) (ii,iv) (Q,S)}
\opt{B}{(I) (ii,iii,iv) (P,S)}
\opt{C}{(II) (iii,iv) (P,R)}
\opt{D}{(IV) (i,iv) (P,R)}
\end{qoptions}
\qn{12} Which of the following options is the only CORRECT combination ?
\begin{qoptions}
\opt{A}{(II) (ii,iii) (P,R)}
\opt{B}{(II) (iii,iv) (Q)}
\opt{C}{(III) (iii,iv) (Q,S)}
\opt{D}{(III) (i,iv) (P,R)}
\end{qoptions}
\qn{13} Which of the following options is the only CORRECT combination ?
\begin{qoptions}
\opt{A}{(III) (i,iii,iv) (P)}
\opt{B}{(III) (ii,iv) (P,Q)}
\opt{C}{(IV) (i,iv) (P,R)}
\opt{D}{(IV) (i,iii,iv) (Q)}
\end{qoptions}
\sect{SECTION-III}
\qn{1} If $y=\left(\ell n\left(x+\sqrt{x^2+1}\right)\right)^2$ and $\left(x^2+1\right)\dfrac{d^3y}{dx^3}+(px+q)\dfrac{d^2y}{dx^2}+r\dfrac{dy}{dx}=sy$, then value of $(p+q+r+s)$ is
\qn{2} Number of points where $f(x)=\left|\left(2^x-1\right).\left(3^x-3\right)x.(x-1).(x-2)\right|$ is not differentiable
\qn{3} Let $f(x)=x+\dfrac{x^2}2+\dfrac{x^3}3+\dfrac{x^4}4+\dfrac{x^5}5$ and $g(x)=f^{-1}(x)$, then $\left|g''(0)\right|$ is
\qn{4} The number of solution of equation $2\sin2x+\sin x(1-2\cos2x)-4=0$ in the interval $(0,\pi)$ :
\qn{5} Let $f(x)=g(x).\mathrm{sgn}\left(x^3+6x^2+11x+6\right)$, where $\mathrm{sgn}$ denotes signum function \& $g(x)=ax^3+x^2+x+1$, then minimum number of points of discontinuity of the function $f(x)$ is
\end{multicols}
\end{document}
